\documentclass[10pt,a4paper]{amsart}
\usepackage[margin=19mm]{geometry}
\usepackage{amsmath,amssymb,mathtools}
\usepackage[scr=rsfs,cal=euler]{mathalfa}
\usepackage{xfrac}
\usepackage[unicode,hidelinks,bookmarks=false]{hyperref}
\newcommand{\ie}{i.e.\ }
\newcommand{\cf}{cf.\ }
\newcommand{\resp}{respectively\ }
\newcommand{\Subsection}{\subsection}
\providecommand{\nobreakdash}{\nobreak\hbox{-}}
\setlength{\emergencystretch}{2em}
\multlinegap0pt
\usepackage[shortlabels]{enumitem}
\usepackage{amssymb}
\usepackage{bbm}
\usepackage{mathtools}
\usepackage{xcolor}
\usepackage{caption}
\usepackage{tikz}
\usepackage{algorithm}\usepackage{algpseudocode}
\newtheorem{lemma}{Lemma}
\def\K{\mathcal{K}}
\title{Efficient spectral bounds on the chromatic number of Hamming, Johnson, and Kneser graph powers}
\author{Finn A. Steinke, Luis M. B. Varona}
\date{Source: arXiv:2601.01962v1 (CC BY 4.0)}
\begin{document}
\maketitle

\noindent\emph{Source and licence.} Efficient spectral bounds on the chromatic number of Hamming, Johnson, and Kneser graph powers. Version arXiv:2601.01962v1 (CC BY 4.0), distributed by arXiv under the Creative Commons Attribution 4.0 International licence, http://creativecommons.org/licenses/by/4.0/, as recorded in the arXiv licence metadata for this exact version and shown on the abstract page https://arxiv.org/abs/2601.01962v1. Full licence notice: \texttt{licenses/arxiv-2601.01962-CC-BY-4.0.txt}.\par\smallskip

\noindent\textit{Editorial scope.} Counted unit: the article's lemma lemma:3.4 (source0.tex lines 431--437). The article's complete original proof of it is delivered with it (source0.tex lines 439--487, one \texttt{proof} environment of 49 lines). No mathematics is written, repaired or re-derived: the statement and the proof are the article's own lines, in the article's own notation, and no retained line is substituted or re-flowed.  No further source-local context is needed: every cross-reference of every retained line resolves inside the two delivered spans.  Nothing was omitted from any retained span, so the package carries no omission record.  No retained line contains a citation, so the wrapper carries no bibliography and no bibliography entry.  No retained line contains a figure environment, an image-inclusion command or drawing code, so no image is delivered and the package declares no asset dependency.  Editorial formatting only, in addition: an AMS \texttt{amsart} preamble that restates the article's own declarations and macros used by the retained lines, namely the environments \texttt{lemma} on separate counters, together with the standard packages those lines need.  The retained environments print in this document as Lemma 1 in order of appearance, whereas the article prints the same items with the numbering of its own full document.  The complete native source of the article as distributed in the arXiv e-print for this exact version is bundled unchanged as \texttt{sources/192125/source0.tex}.
\begin{lemma}\label{lemma:3.4}
    For all $1 \le i \le p$, the following recurrence relations between Kravchuk polynomials hold:
    \begin{enumerate}[label=(\roman*)]
        \item $\displaystyle \K_i(0; n, q) = \frac{(n - i + 1)(q - 1)}{i} \cdot \K_{i - 1}(0; n, q)$
        \item $\displaystyle \K_i(k; n, q) = \frac{(n - k - i + 1)(q - 1)}{i} \cdot \K_{i - 1}(k; n, q) - \frac{k}{i} \cdot \K_{i - 1}(k - 1; n, q)$ for all $k \ge 1$
    \end{enumerate}
\end{lemma}

\begin{proof}
    Fix $n \ge 1$ and $1 \le p \le n$, and suppose that $1 \le i \le p$ and $0 \le k \le n$. Multiplying the Kravchuk polynomial $\K_i(k; n, q)$ by $i$ gives us
    \begin{align*}
        i \cdot \K_i(k; n, q) & = ((i - j) + j) \sum_{j = 0}^i (-1)^j (q - 1)^{i - j} \binom{k}{j} \binom{n - k}{i - j} \\
        & = \sum_{j = 0}^i (-1)^j (q - 1)^{i - j} \binom{k}{j} (i - j) \binom{n - k}{i - j} + \sum_{j = 0}^i (-1)^j (q - 1)^{i - j} j \binom{k}{j} \binom{n - k}{i - j}.
    \end{align*}
    When $j = i$, the expression inside the first sum evaluates to zero, and when $j = 0$, the expression inside the second sum evaluates to zero. Therefore, we can simplify this further to
    \begin{align*}
        i \cdot \K_i(k; n, q) = \sum_{j = 0}^{i - 1} (-1)^j (q - 1)^{i - j} \binom{k}{j} (i - j) \binom{n - k}{i - j} + \sum_{j = 1}^i (-1)^j (q - 1)^{i - j} j \binom{k}{j} \binom{n - k}{i - j}.
    \end{align*}
    We can rearrange the well-known recurrence relation $\binom{a}{b} = \frac{(a - b + 1)}{b}\binom{a}{b - 1}$ to obtain the useful identity $b\binom{a}{b} = (a - b + 1)\binom{a}{b - 1}$. Setting $a = n - k$ and $b = i - j$, we notice that $i \cdot \K_i(k; n, q)$ then becomes
    \begin{align*}
        i \cdot \K_i(k; n, q) ={} & \sum_{j = 0}^{i - 1} (-1)^j (q - 1)^{i - j} \binom{k}{j} (n - k - i + j + 1) \binom{n - k}{i - j - 1} \\
        & + \sum_{j = 1}^i (-1)^j (q - 1)^{i - j} j \binom{k}{j} \binom{n - k}{i - j} \\
        ={} & (n - k - i + 1)(q - 1) \cdot \K_{i - 1}(k; n, q) \\
        & + \sum_{j = 0}^{i - 1} (-1)^j (q - 1)^{i - j} j \binom{k}{j} \binom{n - k}{i - j - 1} + \sum_{j = 1}^i (-1)^j (q - 1)^{i - j} j \binom{k}{j} \binom{n - k}{i - j}.
    \end{align*}
    If $k = 0$, then $\binom{k}{j} = 0$ and the two summations both evaluate to $0$, so dividing both sides by $i$ summarily results in the statement of Lemma \ref{lemma:3.4}(i). With this out of the way, we now consider the case that $k \ge 1$ toward proving Lemma \ref{lemma:3.4}(ii).

    Certainly, when $j = 0$, we have $(-1)^j (q - 1)^{i - j} j \binom{k}{j} \binom{n - k}{i - j - 1} = 0$, so we can simplify further to
    \begin{align*}
        i \cdot \K_i(k; n, q) ={} & (n - k - i + 1)(q - 1) \cdot \K_{i - 1}(k; n, q) \\
        & + \sum_{j = 1}^{i - 1} (-1)^j (q - 1)^{i - j} j \binom{k}{j} \binom{n - k}{i - j - 1} + \sum_{j = 1}^i (-1)^j (q - 1)^{i - j} j \binom{k}{j} \binom{n - k}{i - j}.
    \end{align*}
    By the absorption identity, $j\binom{k}{j} = k\binom{k - 1}{j - 1}$; since we have required $k \ge 1$, this quantity is well-defined. Plugging this in and reindexing with $l = j - 1$ then gives us
    \begin{align*}
        i \cdot \K_i(k; n, q) ={} & (n - k - i + 1)(q - 1) \cdot \K_{i - 1}(k; n, q) \\
        & - k\Biggl(\sum_{l = 0}^{i - 2} (-1)^l (q - 1)^{(i - 1) - l} \binom{k - 1}{l} \binom{n - k}{(i - 1) - l - 1} \\
        & \phantom{\; - k\Biggl(} + \sum_{l = 0}^{i - 1} (-1)^l (q - 1)^{(i - 1) - l} \binom{k - 1}{l} \binom{n - k}{(i - 1) - l}\Biggr).
    \end{align*}
    Isolating the $l = i - 1$ case in the second summation and combining the two resulting summations, we have
    \begin{align*}
        i \cdot \K_i(k; n, q) ={} & (n - k - i + 1)(q - 1) \cdot \K_{i - 1}(k; n, q) \\
        & - k\Biggl(\sum_{l = 0}^{i - 2} (-1)^l (q - 1)^{(i - 1) - l} \binom{k - 1}{l} {\left(\binom{n - k}{(i - 1) - l - 1} + \binom{n - k}{(i - 1) - l}\right)} \\
        & \phantom{\; - k\Biggl(} + (-1)^{i - 1} \binom{k - 1}{i - 1}\Biggr).
    \end{align*}
    We can then apply Pascal's identity $\binom{a}{b} = \binom{a - 1}{b - 1} + \binom{a - 1}{b}$, setting $a = n - k$ and $b = (i - 1) - l$, to obtain
    \begin{align*}
        i \cdot \K_i(k; n, q) ={} & (n - k - i + 1)(q - 1) \cdot \K_{i - 1}(k; n, q) \\
        & - k{\left(\sum_{l = 0}^{i - 2} (-1)^l (q - 1)^{(i - 1) - l} \binom{k - 1}{l} \binom{n - k + 1}{(i - 1) - l} + (-1)^{i - 1} \binom{k - 1}{i - 1}\right)}.
    \end{align*}
    Since both $(q - 1)^{(i - 1) - l}$ and $\binom{n - k + 1}{(i - 1) - l}$ evaluate to $1$ when $l = i - 1$, we can fold the isolated $(-1)^{i - 1} \binom{k - 1}{i - 1}$ term back into the sum to get
    \begin{align*}
        i \cdot \K_i(k; n, q) ={} & (n - k - i + 1)(q - 1) \cdot \K_{i - 1}(k; n, q) \\
        & - k{\left(\sum_{l = 0}^{i - 1} (-1)^l (q - 1)^{(i - 1) - l} \binom{k - 1}{l} \binom{n - k + 1}{(i - 1) - l}\right)} \\
        ={} & (n - k - i + 1)(q - 1) \cdot \K_{i - 1}(k; n, q) - k \cdot \K_{i - 1}({k - 1}; n, q).
    \end{align*}
    Finally, dividing both sides by $i$ yields the statement of Lemma \ref{lemma:3.4}(ii), thus completing the proof.
\end{proof}

\end{document}
